% ============================================================
%  ch02b：§2.2 受限变差的局部系统（原文 pp.9–12）
% ============================================================
\subsection{受限变差的局部系统：导引}

\pnumt{问题何在？}{} 现在假设 $k$ 是代数闭域，$X$ 是 $k$ 上的光滑射影曲线。我们令 $\e := \overline{\Q}_\ell$。

我们希望模仿一般的几何朗兰兹故事，但把局部系统理解为 lisse \et{} $\e$-层，而不是带联络的向量丛。会出什么问题？

首先，让我们尽可能乐观：我们希望有一个叠 $\LS^{\et}_{\bG} = \LS^{\et}_{\bG}(X)$，其行为像先前的叠 $\LS^{\dR}_{\bG}$。在恰当的理解下，它的点应当是 $\bG$-局部系统，即右 $t$-正合对称幺半函子 $\Rep(\bG) \to \qLisse(X)$——这里 $\qLisse(X) \subseteq \Shv(X) = \Shv^{\et}(X)$ 按 \et{} 意义理解，如 §1.4 那样。

作为预备步骤，注意 $X$ 上平凡 $\bG$-局部系统的自同构群是 $\bG(\e)$。这暗示 $\LS^{\et}_{\bG}$ 应定义在域 $\e$ 上，并包含一个对应于平凡局部系统的 $\mathbf{B}\bG$ 的拷贝。

\begin{remark}
这里我们看到几何朗兰兹的代数几何中的一个基本分叉：某些对象，如 $X$、$G$、$\Bun_G$ 等，生活在\emph{基域} $k$ 上，而谱对象，如 $\bG$、$\LS^{\et}_{\bG}$ 等，生活在\emph{系数域} $\e$ 上。对 de Rham 理论，系数域就是基域，这一区分可以忽略。
\end{remark}

\pnum{}现在假设 $k$ 特征为零，且 $X$ 的亏格 $g > 0$。

忽略技术问题（叠性、导出结构），我们可能首先猜测 $\LS^{\et}_{\mathbb{G}_m}$ 是域 $\e$ 上类似 $\mathbb{G}_m^{2g}$ 的东西。毕竟，$X$ 的 \et{} 基本群的 Abel 化是 $\widehat{\ZZ}^{2g}$。此外，人们可以看出（忽略同样的技术问题），叠 $\LS^{\dR}_{\mathbb{G}_m}$ 在 $\mathbb{C}$ 上\emph{复解析地}同构于 $\mathbb{G}_m^{2g} \times (\mathbf{B}\mathbb{G}_m \times 0 \times_{\mathbb{A}^1} 0)$。

然而，$\ZZ$ 与 $\widehat{\ZZ}$ 之间的区别在这里是关键。事实上，连续同态 $\widehat{\ZZ}^{2g} \to \e^\times$ 由 $(\OO^{\times}_{\e})^{2g}$ 中的点参数化，其中 $\OO_{\e} \subseteq \e$ 是整元素的通常赋值子环。换言之，我们希望 $\LS^{\et}_{\mathbb{G}_m} = \mathbb{G}_m^{2g} \times (\mathbf{B}\mathbb{G}_m \times 0 \times_{\mathbb{A}^1} 0) \approx \mathbb{G}_m^{2g}$ 太天真了：右端在 $\e$ 上点太多了！

注意，很难找到一个有趣的\footnote{这里有一个并不有趣的例子：$\prod_{\OO^\times_{\e}} \Spec(\e)$。对 $\mathbb{G}_m$ 的情形，$\LS^{\et}_{\mathbb{G}_m}$ 将在这个答案的基础上，于叠性与（可能的）上同调幂零元方面有所不同（当然，还要把 $\OO^\times_{\e}$ 换成 $(\OO^\times_{\e})^{2g}$）。}、以 $\e$ 为系数域且具有 $\e$-点 $\OO^{\times}_{\e}$ 的（譬如说连通且代数的）叠 $\LS^{\et}_{\mathbb{G}_m}$。所以我们放弃。

\pnumt{另一方面，}{} \et{} 局部系统（别名：Galois 表示）的形变理论是一个老故事。通常人们考虑挠系数，但我们在这里不必如此。要点是：对 \et{} 局部系统 $\sigma$，我们有一个 DG 李代数 $C_{\et}(X, \check{\mathfrak{g}}_\sigma)$，因此有一个相伴的形式模问题（参见 [25] 第 13 章与 [18]）。

换言之，尽管我们放弃了 $\LS^{\et}_{\bG}$ 的存在性，我们\emph{确实}知道它的 $\e$-点（毕竟那都是局部系统），并且\emph{确实}知道它在每个这样的点处的形式完备化。

所以目前，为了构造 $\LS^{\et}_{\bG}$ 的\emph{某个}近似，我们可以取“$\LS^{\et}_{\bG}$ 在每个 $\e$-点 $\sigma$ 处的形式完备化”的（通常是不可数的）无交并。这对 $\mathbb{G}_m$ 给出了正确答案，但对其他群我们将构造一个稍好一些的东西（正如 §2.2 余下部分所讨论的）。

\pnumt{定义。}{} 在 [2] 中，我们对任意仿射代数群 $H/\e$ 定义了 $\e$ 上的一个\emph{预叠}（即从连通交换 $\e$-代数到 $\infty$-群胚的函子，也就是模问题），称为 $\LS^{\restr}_{H}$。一般地，它比点的形式邻域记忆得稍多一点，正如我们将看到的。

这个定义是对 (2.3) 的天真模仿；按定义，$\LS^{\restr}_{H}$ 的一个 $S$-点是右 $t$-正合对称幺半函子：
\[
  \Rep(H) \to \qLisse(X) \otimes \QCoh(S).
\]

\begin{remark}
为使这个定义显得更具体，我们解释一下不用 DG 范畴张量积时右端是什么。假设 $\mathcal{C}$ 是一个 DG 范畴，我们提醒它有所有余极限。假设 $S = \Spec(A)$。于是 $\mathcal{C} \otimes \QCoh(S) = \mathcal{C} \otimes A\text{-}\mathrm{mod} = A\text{-}\mathrm{mod}(\mathcal{C})$，即 $\mathcal{C}$ 中带 $A$-作用的对象。所以右端相当具体——其复杂程度与 $A$-$\mathrm{mod}$ 大致相同。
\end{remark}

\pnumt{我们得到了什么？}{} 上面是一个形式定义。仍需论证我们给出了一个\emph{好的概念}，而“好”的含义将随着我们前进而变得更精细。

\pnumt{例：加性群。}{} 首先假设 $H = \mathbb{G}_a$。

我们断言，在这种情形中 $\LS^{\restr}_{\mathbb{G}_a}$ 实际上“看起来与” de Rham 情形相同！更确切地说，我们将证明 $\LS^{\restr}_{\mathbb{G}_a}$ 是对应于复形 $C_{\et}(X)[1]$ 的代数叠，即它（非典范地\footnote{我们实际上为 $C_{\et}(X)[1]$ 选择了一个形式性同构。}）同构于
\[
  \mathbf{B}H^0_{\et}(X) \times H^1_{\et}(X) \times \Omega H^2_{\et}(X)
\]
（其中 $\Omega H^2_{\et}(X)$ 定义为导出概形 $0 \times_{H^2_{\et}(X)} 0$）。人们可以看出（例如通过下面的分析），在 de Rham 设定中同样成立，只是处处用 de Rham 上同调替换 \et{} 上同调。

为看出这一点，令 $\triv \in \Rep(\mathbb{G}_a)$ 为平凡表示，即张量单位。存在一个典范映射 $\triv \to \triv[1] \in \Rep(\mathbb{G}_a)$，对应于 $\mathbb{G}_a$ 的标准 $2$ 维表示 $\begin{pmatrix} 1 & 1 \\ 0 & 1 \end{pmatrix}$（看作平凡表示的非平凡自扩张）。

现在，对对称幺半 DG 范畴 $\mathcal{C}$ 与对称幺半函子 $F: \Rep(\mathbb{G}_a) \to \mathcal{C}$，我们可以把 $F$ 作用到上面的扩张类上，得到一个映射 $\mathbf{1}_{\mathcal{C}} \to \mathbf{1}_{\mathcal{C}}[1]$，即 $(\infty\text{-})\text{群胚}$ $\Hom_{\mathcal{C}}(\mathbf{1}_{\mathcal{C}}, \mathbf{1}_{\mathcal{C}}[1]) = \Omega^{\infty-1}\underline{\mathrm{End}}_{\mathcal{C}}(\mathbf{1}_{\mathcal{C}})$ 中的一个点。\footnote{为向不熟悉这些材料的读者澄清：对谱 $V$（或 $k$-向量空间复形），$\Omega^{\infty}V$ 意为“取底层的 $\infty$-群胚”——至少对连通谱/链复形，这类似于取 Abel 群的底层集合；一般地，可以把它设想为“过渡到连通覆盖，然后取底层的同伦集合”。在显式的集合论模型中，我们可以取 $\ZZ$-模的链复形 $V^\bullet$，截断得到 $\tau^{\leq 0}(V^\bullet)$，然后在 Dold--Kan 下过渡到相应的单纯 Abel 群（从而单纯集合）。

记号 $\Omega^{\infty-1}(V)$ 只是指 $\Omega^{\infty}(V[1])$。} 容易看出\footnote{也就是说，只需利用 $\Rep(\mathbb{G}_a)$ 与交换代数 $\e \times \e[-1]$ 上的模范畴之间存在标准的对称幺半等价。}，这给出群胚的同构：
\[
  \Hom_{\mathrm{ComAlg}(\DGCat_{\mathrm{cont}})}(\Rep(\mathbb{G}_a), \mathcal{C}) \simeq \Hom_{\mathcal{C}}(\mathbf{1}_{\mathcal{C}}, \mathbf{1}_{\mathcal{C}}[1]).
\]
取 $\mathcal{C} = \qLisse(X) \otimes \QCoh(S)$，我们看到 $\LS^{\restr}_{\mathbb{G}_a}$ 的 $S$-点等于
\[
  \Omega^{\infty}\bigl(\underline{\Hom}_{\qLisse(X)}(\mathbf{e}_X, \mathbf{e}_X[1]) \otimes \underline{\Hom}_{\QCoh(S)}(\OO_X, \OO_X)\bigr) = \Omega^{\infty}\bigl(C_{\et}(X)[1] \otimes \Gamma(S, \OO_S)\bigr),
\]
其中 $\mathbf{e}_X$ 是 $X$ 上的常层。把形式体系展开之后，这就证明了断言。

\pnumt{例：乘性群。}{} 这里我们只陈述结果：

空间 $\LS^{\restr}_{\mathbb{G}_m}$ 是一个 ind-代数叠。它是其连通分支的无交并，而每个连通分支（同样非典范地\footnote{这里有一个更典范地构造该分支的方法。首先取 $\LS^{\restr}_{\mathbb{G}_a}$ 并在平凡 $\mathbb{G}_a$-局部系统处形式完备化。所得的叠接受来自 $\mathbf{B}\mathbb{G}_a^{\wedge} = \mathbf{B}\mathbb{G}_m$ 的同态（这里我们用了指数映射）；然后沿到 $\mathbf{B}\mathbb{G}_m$ 的映射作推出。}）同构于 $\mathbf{B}\mathbb{G}_m \times H^1_{\et}(X)^{\wedge}_0 \times \Omega H^2_{\et}(X)$。$\LS^{\restr}_{\mathbb{G}_m}$ 的连通分支与其 $\e$-点一一对应，我们提醒后者恰是 $X$ 上秩 $1$ 的 lisse 层。

\pnumt{工具箱是什么？}{} 这部分材料可以跳过。为方便读者，我们描述证明关于 $\LS^{\restr}_{H}$ 的事情的常用方法。

首先，我们需要探查底层的经典叠，忽略导出代数几何的问题。为此，令 $\Pi_X$ 为附着于 Tannaka 范畴 $\qLisse(X)^{\heartsuit}$ 的 Tannaka 群，于是 $\Pi_X$ 是 $\e$ 上的群概形，并具有对称幺半等价 $\Rep(\Pi_X)^{\heartsuit} \simeq \qLisse(X)^{\heartsuit}$。

对经典概形，$\LS^{\restr}_{H}$ 的 $S$-点典范地双射于群概形映射 $\Pi_X \times S \to H \times S$（模共轭）（其中我们在群胚意义下取商）——参见 [2] 命题 2.5.9（尽管该断言本质上只是 Tannaka 对偶加上簿记）。这允许我们使用代数群论的工具研究 $\LS^{\restr}_{H}$ 的底层经典预叠。

然后我们用形变理论推广到导出概形，这对 $\LS^{\restr}_{H}$ 来说容易计算；参见 [2] §2.2。

\begin{example}
我们用一个例子说明第一种技术。

先前关于 $\LS^{\restr}_{\mathbb{G}_m}$ 是 §2.2.7 中“胖点”的无交并的断言，等价于说：对任意代数闭域扩张 $\e'/\e$，映射 $S = \Spec(\e') \to \LS^{\restr}_{\mathbb{G}_m}$ 通过一个 $\e$-点分解。这变成关于群概形的一般断言：映射 $\Pi_X \times S \to \mathbb{G}_m \times S$ 来自一个定义在 $\e$ 上的映射。由于 $\mathbb{G}_m$ 具有有限型，这归约到同样的断言，但把 $\Pi_X$ 换成仿射代数群 $\Gamma$（即 $\Pi_X$ 的有限型商），我们甚至可以假设它是 Abel 的。这里该断言由交换代数群的表示论显然可见。
\end{example}

\pnumt{$\LS^{\restr}$ 的一般结构。}{} 我们希望下面的结果在此时不会令人惊讶。

首先，$\LS^{\restr}_{H}$ 总是一个\emph{形式}代数叠。更确切地说，若把 $\LS^{\restr}_{H} = \LS^{\restr}_{H}(X) \to \mathbf{B}H = \LS^{\restr}_{H}(\Spec(k))$ 取为在 $X$ 的一个 $k$-点处的纤维，则该映射在 indscheme 中可表，而且更好：在形式概形的无交并这些 indscheme 中可表——参见 [2] 定理 1.4.5。

其次，$\LS^{\restr}_{H}$ 的连通分支与\emph{半单} $H$-局部系统（模等价）一一对应。非形式地说，$\LS^{\restr}_{H}$ 的两个点属于同一连通分支，当且仅当它们的半单化在无穷小意义下接近。确切陈述参见 [2] 命题 3.7.2。

最后，如果我们设想 $\LS^{\et}_{H}$ 存在，那么对每个半单 $\sigma$，将存在局部系统的一个闭子叠 $\LS^{\et}_{H,\sigma}$，其半单化为 $\sigma$；此时 $\LS^{\restr}_{H}$ 在道德上就是 $\LS^{\et}_{H}$ 在每个这样的 $\LS^{\et}_{H,\sigma}$ 处形式完备化的无交并。关于 Betti 与 de Rham 设定中更精确的断言，参见 [2] §4。
